\documentclass{mcmthesis}
\mcmsetup{CTeX = false,                  % 如果只写英文，设为false；如需中文注释或内容，设为true
          tcn = {2617170},               % 【修改此处】把这里改成你们的队伍控制号 Team Control Number
          problem = \textcolor{red}{D},  % 【修改此处】选择的题号 A/B/C/D/E/F
          sheet = true,                  % 是否输出摘要页（Summary Sheet），美赛必须为true
          titleinsheet = true, 
          keywordsinsheet = true,
          titlepage = false,             % 美赛通常不需要单独的标题页
          abstract = true,
          }

% ==================== 常用宏包引入 ====================
\usepackage{palatino}     % 使用 Palatino 字体，看起来比 Times 更现代、优雅
\usepackage{amsmath, amssymb, amsfonts}  % 数学公式与符号
\usepackage{booktabs}     % 创建专业的三线表
\usepackage{longtable}    % 跨页长表格
\usepackage{graphicx}     % 插入图片
\usepackage{subcaption}   % 并排子图
\usepackage{float}        % 强制图片位置 [H]
\usepackage{listings}     % 代码块高亮
\usepackage{xcolor}       % 颜色支持
\usepackage{geometry}     % 页面布局
\usepackage{lipsum}       % 生成填充文本（实际写作时删除）
\usepackage{comment}      % 可以进行注释          
\usepackage{newtxtext}    % \usepackage{palatino}
\usepackage[backend=bibtex]{biblatex}   % for RStudio Complie
\usepackage{tocloft}

% ==================== 代码块设置 ====================
% 针对 Python/C++ 优化的代码样式
\definecolor{codegreen}{rgb}{0,0.6,0}
\definecolor{codegray}{rgb}{0.5,0.5,0.5}
\definecolor{codepurple}{rgb}{0.58,0,0.82}
\definecolor{backcolour}{rgb}{0.95,0.95,0.92}

\lstdefinestyle{mystyle}{
    backgroundcolor=\color{backcolour},   
    commentstyle=\color{codegreen},
    keywordstyle=\color{magenta},
    numberstyle=\tiny\color{codegray},
    stringstyle=\color{codepurple},
    basicstyle=\ttfamily\footnotesize,
    breakatwhitespace=false,         
    breaklines=true,                 
    captionpos=b,                    
    keepspaces=true,                 
    numbers=left,                    
    numbersep=5pt,                  
    showspaces=false,                
    showstringspaces=false,
    showtabs=false,                  
    tabsize=2
}
\lstset{style=mystyle}
\setlength{\cftbeforesecskip}{6pt}
\renewcommand{\contentsname}{\hspace*{\fill}\Large\bfseries Contents \hspace*{\fill}}

\title{Enjoy a Cozy and Green Bath}
% \author{\small \href{http://www.latexstudio.net/}
%   {
\includegraphics[width=7cm]{mcmthesis-logo}}}
\date{\today}

\begin{document}

\begin{abstract}

A traditional bathtub cannot be reheated by itself, so users have to add hot water from time to time. Our goal is to establish a model of the temperature of bath water in space and time. Then we are expected to propose an optimal strategy for users to keep the temperature even and close to initial temperature and decrease water consumption.

To simplify modeling process, we firstly assume there is no person in the bathtub. We regard the whole bathtub as a thermodynamic system and introduce heat transfer formulas.

We establish two sub-models: adding water constantly and discontinuously. As for the former sub-model, we define the mean temperature of bath water. Introducing Newton cooling formula, we determine the heat transfer capacity. After deriving the value of parameters, we deduce formulas to derive results and simulate the change of temperature field via CFD. As for the second sub-model, we define an iteration consisting of two process: heating and standby. According to energy conservation law, we obtain the relationship of time and total heat dissipating capacity. Then we determine the mass flow and the time of adding hot water. We also use CFD to simulate the temperature field in second sub-model.

In consideration of evaporation, we correct the results of sub-models referring to some scientists' studies. We define two evaluation criteria and compare the two sub-models. Adding water constantly is found to keep the temperature of bath water even and avoid wasting too much water, so it is recommended by us.

Then we determine the influence of some factors: radiation heat transfer, the shape and volume of the tub, the shape/volume/temperature/motions of the person, the bubbles made from bubble bath additives. We focus on the influence of those factors to heat transfer and then conduct sensitivity analysis. The results indicate smaller bathtub with less surface area, lighter personal mass, less motions and more bubbles will decrease heat transfer and save water.

Based on our model analysis and conclusions, we propose the optimal strategy for the user in a bathtub and explain the reason of uneven temperature throughout the bathtub. In addition, we make improvement for applying our model in real life.

\begin{keywords}
Heat transfer, Thermodynamic system, CFD, Energy conservation
\end{keywords}

\end{abstract}

\maketitle

%% Generate the Table of Contents, if it's needed.
% \renewcommand{\contentsname}{\centering Contents}
\tableofcontents        % 若不想要目录, 注释掉该句
\thispagestyle{empty}

\newpage

\section{Introduction}


\subsection{Problem Background}
Professional sports represent a unique intersection of athletic competition and business enterprise. The WNBA provides a particularly relevant setting for studying this trade-off due to its rapid growth in media exposure, increasing franchise valuations, and strict structural constraints such as salary caps and roster limits.

While the framework we develop is applicable across WNBA franchises, we focus on the Indiana Fever as a representative case study.The Fever offers a uniquely suitable analytical setting due to the concentration of commercial value in a single superstar player (Caitlin Clark), recent revenue expansion, and publicly observable performance and financial indicators.

\subsection{Problem Restatement}
To address the objectives of this problem, we must construct an integrated analytical framework that supports strategic decision-making for a professional WNBA team. Specifically, we aim to:

\begin{itemize}
\item {\bf Quantifying individual player value} We will quantify their value by incorporating both on-court performance and commercial impact.

\item {\bf Developing a financial model} The model will link team composition, operating costs, and revenue generation to overall profitability and franchise value.
\item {\bf Establishing a dynamic decision-making mechanism} We'll adjust roster and investment strategies in response to changes in performance, injury status, and economic conditions.

\item {\bf Evaluating player acquisition strategies} We'll discuss strategies under league constraints such as salary caps and roster limits.

\item {\bf Assessing the impact of external structural changes} We'll put typical changes into consideration,such as league expansion, on team profitability and competitive positioning.
\end{itemize}

\subsection{Team selection}
\subsection{Problem Analysis}
The problem contains several interconnected components that require structured modeling.

First, player evaluation must extend beyond traditional performance statistics. Athletic contribution influences team wins, while popularity and market appeal affect revenue streams such as ticket sales, merchandise, and media exposure. Therefore, player value should be modeled as a multi-dimensional construct combining competitive and commercial factors.

Second, team profitability depends on both revenue generation and cost control. Revenue is influenced by team performance, player marketability, and external market conditions, whereas costs are constrained by salary structures and league regulations. A financial optimization framework is required to determine sustainable investment levels.

Third, decision-making is dynamic rather than static. Player injuries, aging effects, and economic fluctuations introduce uncertainty over time. Consequently, the model must incorporate risk and allow for strategic adjustment across multiple periods.

Finally, league-level policies, including salary caps and potential expansion, impose structural constraints that alter optimal strategies. These constraints must be explicitly embedded within the modeling framework.

\section{Assumptions and Justifications}

To simplify the problem and make it convenient for us to simulate real-life conditions, we make the following basic assumptions, each of which is properly justified.

\begin{itemize}
\item {\bf Player performance can be represented by aggregated statistical indicators.}

We assume that a player’s athletic contribution can be sufficiently captured by key performance metrics such as scoring efficiency, rebounds, assists, and defensive impact, which are publicly available and measurable.
We further assume that Wins Above Replacement (WAR), derived from Box Plus/Minus (BPM), provides a valid and consistent proxy for marginal contribution to team wins.

Aggregated statistics provide an objective and standardized representation of on-court contribution. These indicators are widely used in professional analytics and allow for quantitative comparison across players.

\item {\bf Commercial impact of a player can be approximated using measurable proxies.}

We assume that a player’s commercial value can be estimated using observable indicators such as media exposure, popularity, and merchandise influence.

Direct financial contribution of individual players is difficult to isolate due to shared team revenue streams.So,proxy variables are necessary to approximate market influence in a consistent and data-driven manner.

\item {\bf Team revenue is positively correlated with competitive performance and player marketability.}

We assume that improved team performance increases revenue streams such as ticket sales, sponsorship value, and media exposure. Additionally, high-profile players contribute independently to revenue generation.

Empirical evidence from professional sports economics suggests that winning teams attract higher attendance and commercial attention. Furthermore, superstar players often generate revenue beyond performance-related outcomes.

\item {\bf Salary cap and roster constraints are binding and strictly enforced.}

We assume that league regulations, including salary caps and roster limits, restrict feasible roster construction strategies.

WNBA teams operate under formal league governance structures that limit payroll flexibility. These constraints must be embedded in the optimization framework to ensure realistic and implementable solutions.

\item {\bf  Player aging and injury risk affect future performance probabilistically.}

We assume that player performance declines gradually with age and may be reduced under injury scenarios, modeled through risk-adjusted performance factors.

Athletic performance follows a life-cycle pattern in professional sports. Incorporating age and injury uncertainty improves the realism of long-term value estimation and supports dynamic decision-making.

\item{\bf Management decisions are rational and profit-oriented.}

We assume that team management aims to maximize a weighted objective function that balances competitive success and financial return.

Although real-world decisions may include behavioral factors, modeling rational optimization allows us to derive systematic and interpretable strategies aligned with ownership objectives.

\end{itemize}

\section{Notations}

\begin{table}[h]
\centering
\caption{Notation Summary}
\begin{tabular}{l l c}
\toprule
\textbf{Symbol} & \textbf{Description} & \textbf{Unit} \\
\midrule
$V_{total}$ & Total player value & USD \\
$V_{perf}$ & Performance-based value & USD \\
$V_{comm}$ & Commercial value & USD \\
$\beta$ & Weight assigned to performance value & -- \\
$WAR$ & Wins Above Replacement & wins \\
$BPM$ & Box Plus/Minus & points/100 poss \\
$VPW$ & Value Per Win & USD/win \\
$R$ & Total revenue & USD \\
$C$ & Total costs & USD \\
$\pi$ & Operating profit & USD \\
$Cap$ & Salary cap & USD \\
$WACC$ & Weighted Average Cost of Capital & \% \\
$D/V$ & Debt-to-Value ratio & -- \\
$Overlap$ & Market overlap coefficient & -- \\
$d_{ij}$ & Distance between cities $i$ and $j$ & km \\
\bottomrule
\end{tabular}
\end{table}

%\noindent where we define the main parameters while specific value of those parameters will be given later.

\section{Model Overview}

The overall framework consists of four interconnected modules:

(1) a dual-dimension player valuation model integrating performance and commercial impact; 

(2) a financial forecasting and capital structure model linking roster composition to profitability; 

(3) a mixed-integer roster optimization program under salary cap constraints; 

(4) a trigger-based dynamic decision engine that adjusts strategies in response to injuries, performance thresholds, and structural shocks. 

These modules operate iteratively across the season.

%\begin{figure}[h] 
%\centering
%\includegraphics[width=12cm]{fig1.jpg}
%\caption{Modeling process} \label{fig1}
%\end{figure}

\section{Dynamic Decision-Making Model}

we construct a dynamic decision-making model that integrates basketball operations and financial management within a unified and adaptive structure.
We formulate team management as a multi-period optimization problem in which ownership seeks to balance competitive success and financial sustainability. The model incorporates performance dynamics, revenue fluctuations, and structural constraints to support rational adjustment of roster and investment strategies over time.

\subsection{Dual-Track Management Framework}

We model team governance as two parallel yet interdependent management tracks.



\subsubsection{Business Operations Track}

The Business Operations Track focuses on revenue maximization, cost control, capital allocation, and long-term franchise valuation. Key decision variables include pricing strategy, sponsorship contracts, investment in facilities, and capital structure choices.

Financial performance is evaluated through operating profit:
\[
\Pi = R - C,
\]
where revenue $R$ depends on team performance and market appeal, and cost $C$ includes salaries, operating expenses, and capital costs.
%rectangular coordinate system. The length, width, and height of it is $a,\, b$ and $c$.

\begin{figure}[h] 
\centering
\includegraphics[width=8cm]{fig2.jpg}
\caption{Modeling process} \label{fig2}
\end{figure}

%In the basis of this, we introduce the following equations \cite{5}:

\begin{itemize}
\item {\bf Continuity equation:}
\end{itemize}

\begin{equation} \label{eq1}
\frac{\partial u}{\partial x} + \frac{\partial v}{\partial y} +\frac{\partial w}{\partial z} =0
\end{equation}

%\noindent where the first component is the change of fluid mass along the $X$-ray. The second component is the change of fluid mass along the $Y$-ray. And the third component is the change of fluid mass along the $Z$-ray. The sum of the change in mass along those three directions is zero.

\begin{itemize}
\item {\bf Moment differential equation (N-S equations):}
\end{itemize}



\begin{itemize}
\item {\bf Energy differential equation:}
\end{itemize}

\begin{equation} \label{eq3}
\rho c_p \Big( u\frac{\partial t}{\partial x} + v\frac{\partial t}{\partial y} + w\frac{\partial t}{\partial z} \Big) = \lambda \Big(\frac{\partial^2 t}{\partial x^2} + \frac{\partial^2 t}{\partial y^2} + \frac{\partial^2 t}{\partial z^2} \Big)
\end{equation}

\noindent where the left three components are convection terms while the right three components are conduction terms.

By Equation \eqref{eq3}, we have ......

......

On the right surface in Fig. \ref{fig2}, the water also transfers heat firstly with bathtub inner surfaces and then the heat comes into air. The boundary condition here is ......

\subsubsection{Basketball Operations Track}

The Basketball Operations Track concentrates on competitive success. Its objectives include win maximization, roster optimization, player development, and risk management related to injuries and aging.

Team performance is summarized through a competitive index $W$, which may represent total wins, playoff probability, or a weighted performance metric derived from player contributions.

\subsubsection{Integrated Objective Function}

Although the two tracks pursue distinct operational goals, they are strategically interconnected. Strong competitive performance enhances revenue streams, while financial strength supports talent acquisition and infrastructure investment.

To formalize this relationship, we define a unified objective function:

\[
\max V = \alpha \cdot \Pi + (1-\alpha) \cdot W,
\]

where $\alpha \in [0,1]$ represents the relative weight assigned to financial return versus competitive achievement.

The parameter $\alpha$ captures ownership preference and strategic phase. For example, rebuilding periods may require higher emphasis on financial stability, while championship windows justify greater emphasis on competitive performance.

\subsection{Core Model Components}

To operationalize the dual-track decision framework, we construct three interconnected quantitative sub-models: a Player Value Assessment Model, a Financial Forecasting Model, and a Capital Structure Optimization Model. These components collectively translate strategic objectives into measurable decision variables.

\subsubsection{Player Value Evaluation Sub-Model}

Central to the system is a dual-dimension player valuation model capturing both competitive contribution and commercial influence. The total value of player $i$ is defined as

\[
V_{total,i} = \beta \cdot V_{performance,i} + (1-\beta) \cdot V_{commercial,i},
\]

where $\beta \in [0,1]$ reflects the relative importance of athletic performance.

\textbf{Performance Value.}  
On-court contribution is measured using Wins Above Replacement (WAR), which estimates a player's marginal impact relative to a replacement-level athlete:

\[
WAR_i = \frac{(BPM_i - BPM_{replacement}) \times Minutes_i}{48 \times G},
\]

where $BPM_{replacement}$ represents replacement-level performance and $G$ denotes season length.

The monetary conversion is given by:

\[
V_{performance,i} = WAR_i \times VPW,
\]

where $VPW$ denotes the estimated economic value per additional win.

\textbf{Commercial Value.}  
Commercial contribution reflects revenue generation beyond performance metrics. Since individual revenue attribution is difficult to isolate, we introduce a normalized market influence index $M_i$ incorporating observable indicators such as media exposure and merchandise impact.

Inspired by cooperative game theory, commercial value is allocated proportionally:

\[
V_{commercial,i} = \frac{M_i}{\sum_j M_j} \cdot (R_{total} - R_{baseline}),
\]

where $R_{baseline}$ represents revenue independent of star effects.

This formulation allows objective comparison between players with different competitive and market profiles.

\subsubsection{Financial Forecasting Sub-Model}

We model team profitability as:

\[
\Pi = R_{total} - C_{total},
\]

where total revenue is decomposed into multiple streams:

\[
R_{total} = R_{tickets} + R_{TV} + R_{sponsorship} + R_{merchandise} + R_{other}.
\]

Ticket revenue follows a dynamic pricing mechanism:

\[
R_{tickets} = \sum_{g=1}^{G} P_g \times A_g,
\]

where game-level ticket price $P_g$ adjusts according to opponent strength, scheduling effects, and star availability.

Costs are constrained by salary cap regulations:

\[
\sum_{i=1}^{n} Salary_i \leq Cap.
\]

This ensures financial feasibility and compliance with league governance.

\subsubsection{Capital Structure Optimization}

Beyond operational profit, franchise valuation depends on financing decisions. We minimize the weighted average cost of capital (WACC):

\[
WACC = \frac{E}{V} r_e + \frac{D}{V} r_d (1 - T_c),
\]

where $E$ and $D$ denote equity and debt, respectively.

To reflect increasing financial risk at higher leverage levels, we model the cost of debt as:

\[
r_d = r_{base} + \gamma \left( \frac{D}{V} \right)^2,
\]

where $\gamma$ captures nonlinear risk premium effects.

The optimal capital structure balances financial flexibility, risk exposure, and long-term franchise appreciation.


\subsection{Dynamic Decision Engine}

To ensure adaptability under uncertainty, we introduce a trigger-based adjustment mechanism. The system continuously monitors key state variables across four dimensions: competitive performance, revenue growth, player health, and regulatory environment.

When predefined threshold conditions are met, the model updates strategic parameters, including the objective weight $\alpha$, investment intensity, and roster adjustment policies.

Formally, we define a state vector:

\[
S_t = (W_t, R_t, Injury_t, Policy_t),
\]

where each component represents a measurable condition at time $t$.

The decision parameter $\alpha$ is updated according to:

\[
\alpha_{t+1} = \alpha_t + \delta(S_t),
\]

where $\delta(\cdot)$ represents an adjustment function triggered by threshold conditions.

For example:
\begin{itemize}
\item Sustained revenue growth may increase financial prioritization;
\item Declining win percentage may shift emphasis toward competitive investment;
\item Major injury events may activate short-term roster reinforcement strategies.
\end{itemize}

This mechanism transforms the framework into a responsive decision engine capable of adapting to different strategic phases.

%\subsubsection{Trigger Conditions}
%\subsubsection{Weight Adjustment Mechanism}
%\subsubsection{Risk Control Mechanism}

\subsection{Implementation and Adjustment Path}

We divide strategic execution into three temporal horizons:

\begin{itemize}
\item \textbf{Short-term stabilization:} ensure salary cap compliance and optimize immediate roster efficiency.
\item \textbf{Mid-term optimization:} adjust roster composition based on updated valuation metrics and revenue forecasts.
\item \textbf{Long-term value maximization:} align capital structure and competitive planning with franchise appreciation goals.
\end{itemize}

This staged implementation ensures that dynamic adjustments remain consistent with long-term strategic objectives.
%\subsubsection{Short-Term Strategy}
%\subsubsection{Mid-Term Optimization}
%\subsubsection{Long-Term Franchise Value Maximization}


\section{Player Acquisition Strategy}

Building upon the dynamic decision-making framework developed in Section 3, we formulate player acquisition as a constrained optimization problem. The objective is to construct a roster that maximizes total team value while satisfying financial and regulatory constraints.

\subsection{Roster Optimization Model}

We formulate roster construction as a mixed-integer linear programming (MILP) problem. Let $x_i \in \{0,1\}$ indicate whether player $i$ is selected, $V_i$ denote the total player value, and $S_i$ represent salary.

The objective is:

\[
\max \sum_i V_i x_i
\]

subject to:

\[
\sum_i S_i x_i \leq Cap,
\]

\[
n_{min} \leq \sum_i x_i \leq n_{max},
\]

\[
\sum_{i \in Pos_k} x_i \geq L_k, \quad \forall k,
\]

where $Cap$ is the salary cap, $n_{min}, n_{max}$ denote roster size limits, and $L_k$ represents positional requirements.

For teams with established franchise cornerstones, we define a core player set $\mathcal{C}$ such that:

\[
x_i = 1, \quad \forall i \in \mathcal{C}.
\]

This formulation ensures competitive balance, financial feasibility, and structural stability.


\subsection{Player Potential Projection}

To support multi-year contract decisions, we project future player value using an age-adjusted performance function:

\[
V_{projected}(t) = V_{current} \cdot f(age_t) \cdot (1 - \theta_t),
\]

where $f(age)$ captures position-specific performance trajectories, and $\theta_t$ represents injury-adjusted risk.

Contract valuation is computed through Net Present Value (NPV):

\[
NPV = \sum_{t=1}^{T} \frac{V_{projected}(t)}{(1+r)^t},
\]

where $r$ is the discount rate reflecting franchise capital cost.

This framework prevents overvaluation of aging or high-risk players and promotes sustainable investment.


\subsection{Strategic Acquisition Priorities}

Based on the optimization results, roster enhancement strategies generally fall into three categories:

\begin{itemize}
\item \textbf{Interior Defensive Reinforcement:} Improve defensive efficiency through acquisition of rim-protecting players with positive defensive impact metrics.

\item \textbf{Backcourt Depth Enhancement:} Ensure offensive continuity and load management by adding secondary ball handlers.

\item \textbf{Draft Capital Optimization:} Allocate draft resources toward high-upside players with positional versatility and favorable cost-to-value ratios.
\end{itemize}

These priorities are dynamically adjusted according to the parameter $\beta$ and the strategic phase defined in Section 3.

\subsection{Trade Evaluation Framework}

Potential trades are evaluated using value differential analysis:

\[
\Delta V = V_{incoming} - V_{outgoing}.
\]

A trade is recommended when:

\[
\Delta V > 0
\]

and all salary cap and roster constraints remain satisfied.

Beyond quantitative evaluation, strategic fit is assessed through positional alignment and timeline consistency within the competitive window.



\section{League Expansion Impact Analysis}

League expansion alters competitive dynamics by redistributing market attention, sponsorship allocation, and media exposure. To evaluate these effects, we construct a geographic competition model based on market overlap and distance decay principles.
%\quad~ We first give the value of parameters based on others’ studies. Then we get the calculation results and simulating results via those data.

\subsection{Geographic Market Competition Model}

We define the market overlap between franchise $i$ and franchise $j$ as:

\[
Overlap_{ij} = \frac{\min(Pop_i, Pop_j)}{Pop_i} \cdot e^{-\gamma d_{ij}},
\]

where $Pop_i$ and $Pop_j$ denote metropolitan population sizes, $d_{ij}$ represents geographic distance, and $\gamma$ is a distance decay parameter.

The exponential term reflects declining fan substitution probability with increasing distance. This structure captures the intuition that nearby franchises compete more intensely for regional attention and sponsorship.

\subsection{Revenue Impact Estimation}

We estimate revenue impact as a function of market overlap:

\[
\Delta R_i = - \lambda \sum_{j \in Expansion} Overlap_{ij} \cdot R_i,
\]

where $\lambda$ represents the sensitivity of revenue to overlapping competition.

Expansion locations are classified into four competitive tiers based on geographic distance:

\begin{itemize}
\item \textbf{Critical Threat:} $d < d_1$
\item \textbf{High Threat:} $d_1 \leq d < d_2$
\item \textbf{Moderate Threat:} $d_2 \leq d < d_3$
\item \textbf{Low Threat:} $d \geq d_3$
\end{itemize}

Close-proximity expansion generates direct substitution effects in ticket demand and regional sponsorship allocation, while distant expansion primarily influences league-wide revenue redistribution.

Putting the above value of parameters into the equations we derived before, we can get the some data as follows:

%%普通表格
%\begin{table}[h]  %h表示固定在当前位置
%\centering        %设置居中
%\caption{The calculating results}  %表标题
%\vspace{0.15cm}
%\label{tab2}                       %设置表的引用标签
%\begin{tabular}{|c|c|c|}  %3个c表示3列, |可选, 表示绘制各列间的竖线
%\hline                    %画横线
%Variables & Values & Unit     \\ \hline  %各列间用&隔开
%$A_1$     & 1.05   &   $m^2$  \\ \hline
%$A_2$     & 2.24   &   $m^2$  \\ \hline

%\end{tabular}
%\end{table}

From Table \ref{tab2}, ......

......

\subsection{Strategic Response Under Expansion Scenarios}

After establishing two basic sub-models, we have to correct them in consideration of evaporation heat transfer. Then we define two evaluation criteria to compare the two sub-models in order to determine the optimal bath strategy.

\section{Dynamic Ticket Pricing Strategy}

To enhance revenue performance while maintaining long-term fan engagement, we construct a dynamic ticket pricing model that adapts to demand fluctuations across games.

\subsection{Demand Function Estimation}

We model ticket demand using a linear inverse demand structure:

\[
Q(P) = a - bP,
\]

where $a$ represents maximum potential demand and $b$ captures price sensitivity.

Revenue is given by:

\[
R(P) = P \cdot Q(P) = P(a - bP).
\]

Maximizing revenue yields:

\[
P^* = \frac{a}{2b}.
\]

However, in professional sports settings, stadium capacity $K$ imposes a binding constraint:

\[
Q(P) \leq K.
\]

When unconstrained optimal demand exceeds capacity, the effective price must satisfy:

\[
P \geq \frac{a - K}{b}.
\]

This introduces a capacity-clearing pricing mechanism.



\subsection{Multi-Factor Dynamic Pricing Model}

To account for heterogeneous game characteristics, we introduce demand shifters through multiplicative adjustment factors:

\[
P_g^* = P_{base} \cdot m_{opponent} \cdot m_{schedule} \cdot m_{star} \cdot m_{urgency}.
\]

Each multiplier captures a specific demand driver:

\begin{itemize}
\item Opponent strength effect,
\item Day-of-week scheduling effect,
\item Star availability effect,
\item Purchase timing urgency.
\end{itemize}

These multipliers shift effective demand without altering the underlying price elasticity parameter $b$.
%%三线表
%\begin{table}[h] %h表示固定在当前位置
%\centering  %设置居中
%\caption{Variation of some parameters}  %表标题
%\label{tab7} %设置表的引用标签
%\begin{tabular}{ccccccc} %7个c表示7列, c表示每列居中对齐, 还有l和r可选
%\toprule  %画顶端横线
%$V$      & $A_1$   & $A_2$   & $T_2$    & $q_{m1}$ & $q_{m2}$ & $\Phi_q$ \\
%\midrule  %画中间横线
%-15.00\% & -5.06\% & -9.31\% & -12.67\% & -2.67\%  & -14.14\% & -5.80\% \\
%-12.00\% & -4.04\% & -7.43\% & -10.09\% & -2.13\%  & -11.31\% & -4.63\% \\
%-8.00\%  & -2.68\% & -4.94\% & -6.68\%  & -1.41\%  & -7.54\%  & -3.07\% \\
%-8.00\%  & -2.68\% & -4.94\% & -6.68\%  & -1.41\%  & -7.54\%  & -3.07\% \\
%-8.00\%  & -2.68\% & -4.94\% & -6.68\%  & -1.41\%  & -7.54\%  & -3.07\% \\
%\bottomrule  %画底部横线
%\end{tabular}
%\end{table}

\subsection{Revenue Optimization Under Dynamic Pricing}

Seasonal ticket revenue is computed as:

\[
R_{season} = \sum_{g=1}^{G} P_g^* \cdot \min(Q_g(P_g^*), K).
\]

Compared to static pricing $P_{static}$, dynamic pricing increases revenue when demand variance across games is significant.

To preserve long-term fan loyalty, we impose pricing bounds:

\[
P_{min} \leq P_g^* \leq P_{max}.
\]

This ensures accessibility while preventing negative brand perception from excessive price volatility.




\section{Key Player Injury Management Model}

This section addresses risk adaptation under the sudden loss of a high-value player. 
Professional sports teams often exhibit asymmetric value concentration, where a single player contributes disproportionately to competitive and commercial outcomes. Our framework integrates injury probability modeling, revenue-at-risk estimation, and dynamic response optimization.

\subsection{Injury Probability Model}

We model season-level injury probability as:

\[
P_i = P_0 + \alpha_1 \max(0, Age_i - A_{threshold})
+ \alpha_2 N_{prior}
+ \alpha_3 \max(0, MPG_i - MPG_{threshold})
\]

where:

\begin{itemize}
\item $P_0$ = baseline injury probability,
\item $Age_i$ = player age,
\item $N_{prior}$ = number of prior major injuries,
\item $MPG_i$ = average minutes per game.
\end{itemize}

Expected games missed is modeled as:

\[
E(G_i) = P_i \cdot D_i
\]

where $D_i$ is average injury duration.

\subsection{Dynamic Injury Response Protocol}

We define injury severity levels based on expected duration $D_i$:

\[
Tier =
\begin{cases}
1, & D_i < 3 \text{ weeks} \\
2, & 3 \le D_i < 6 \text{ weeks} \\
3, & D_i \ge 6 \text{ weeks}
\end{cases}
\]

Each tier activates a predefined optimization module:

\begin{itemize}
\item Tier 1: Rotation adjustment
\item Tier 2: Short-term acquisition + marketing reallocation
\item Tier 3: Insurance activation + strategic objective revision
\end{itemize}


%\begin{thebibliography}{99}
%\addcontentsline{toc}{section}{Reference}

%%\bibitem{2} \url{https://en.wikipedia.org/wiki/Convective_heat_transfer#Newton.27s_law_of_cooling}
%\bibitem{3} \url{https://en.wikipedia.org/wiki/Navier\%E2\%80\%93Stokes_equations}
%\bibitem{4} \url{https://en.wikipedia.org/wiki/Computational_fluid_dynamics}
%bibitem{6} Liu Weiguo, Chen Zhaoping, ZhangYin. Matlab program design and application. Beijing: Higher education press, 2002. (In Chinese)

%\end{thebibliography}

%\newpage

\subsection{Load Management Optimization}

We model player workload $W$ as minutes per game. 
Injury probability increases convexly with excessive workload:

\[
P_i(W) = P_0 + \beta (W - W^*)^2
\]

The optimal workload solves:

\[
\min_W \left[ P_i(W) \cdot Loss_i - Marginal\_Win\_Value(W) \right]
\]

This yields an interior solution $W^{opt} < W_{max}$, 
supporting proactive load management policies.

\subsection{Insurance Hedging Decision}

Let premium = $\pi$, coverage = $C$.

Insurance is optimal if:

\[
\pi < P_i \cdot C
\]

For high-value players, this inequality typically holds, 
making disability insurance a rational hedging instrument 
within the broader capital structure strategy.



\section{Sensitivity Analysis}

To evaluate the robustness of our integrated decision-making framework, we conduct sensitivity analysis on key uncertain parameters that materially influence financial and competitive outcomes. Rather than examining isolated components, we focus on system-level stability by perturbing critical inputs within realistic ranges and observing their impact on profitability, roster strategy, and leverage decisions.

The most influential parameter is star player availability. Given the significant commercial weight of core players, particularly Caitlin Clark, a 10\% reduction in games played leads to an approximate \$3.5M decrease in annual profit. This effect propagates through multiple channels, including ticket demand multipliers, merchandise sales, and sponsor exposure. In extreme scenarios involving prolonged injury (6+ weeks), profit reductions can exceed \$10M, validating our emphasis on load management and insurance hedging.

Attendance rate is the second most sensitive variable. Each percentage-point deviation in arena utilization alters annual revenue by approximately \$1.2M. However, dynamic pricing partially mitigates this volatility by adjusting price multipliers in response to demand shifts. As a result, even under moderate demand shocks, revenue fluctuations remain within controllable bounds.

Sponsor retention also exhibits notable influence. Loss of a major sponsor reduces revenue by \$2–3M annually, which in turn affects optimal leverage and investment pacing. Nevertheless, the capital structure recommendation (target debt ratio 20–25\%) remains stable under ±15\% variation in sponsorship-based revenue, indicating resilience in financing strategy.

Conversely, several parameters demonstrate limited structural impact. Variations of ±20\% in the Value Per Win (VPW) estimate do not alter roster optimization outcomes, as relative player rankings remain consistent. Similarly, adjustments to the geographic decay parameter in the expansion model produce only marginal changes in projected revenue impact (<1.5\%), leaving strategic responses unchanged. Draft projection uncertainty also affects long-term value marginally but does not shift short-term acquisition priorities.

Overall, across tested perturbations, the model preserves its core strategic recommendations: maintaining financial-priority weighting (α ≈ 0.55), retaining core players, targeting interior defensive reinforcement, and adopting conservative leverage under CBA uncertainty. This indicates that while short-term profitability is sensitive to star availability and attendance elasticity, the structural decision framework remains robust.

Therefore, the proposed system demonstrates both responsiveness to critical risk factors and stability in long-term strategic orientation.


\section{Model Strengths and Limitations}

\subsection{Strengths}

First, the proposed framework integrates basketball operations and financial management into a unified optimization structure. By combining competitive performance metrics (WAR, BPM) with commercial valuation and capital structure considerations, the model avoids treating athletic and economic decisions independently. This dual-track coordination reflects the real governance structure of professional sports franchises.

Second, the model is dynamically adaptive rather than static. The trigger-based decision engine allows management to respond systematically to performance fluctuations, injury events, expansion threats, and macro-structural changes such as CBA revisions. This enhances practical applicability and ensures that recommendations remain relevant across different states of uncertainty.

Third, the framework incorporates multiple constraints—including salary cap limits, roster composition requirements, leverage boundaries, and market overlap effects—thereby ensuring feasibility of recommended strategies. The use of mixed-integer programming for roster construction further guarantees mathematically consistent and implementable solutions.

Finally, sensitivity analysis demonstrates structural robustness. While profitability is naturally sensitive to star availability and attendance elasticity, core strategic recommendations remain stable under realistic perturbations of key parameters.

\subsection{Limitations}

Despite its strengths, the model relies on several simplifying assumptions. Player performance is represented through aggregated statistics, which may not fully capture qualitative factors such as leadership, locker-room chemistry, or defensive communication. Similarly, commercial value is estimated through proxy indicators rather than direct revenue attribution, potentially introducing estimation bias.

The financial forecasting component assumes relatively stable macroeconomic conditions within a season. Unexpected shocks—such as labor disputes, sudden regulatory changes, or economic downturns—may alter revenue streams beyond the model’s predictive scope.

In addition, the geographic competition model simplifies market behavior through distance-based decay, without incorporating detailed media market segmentation or cross-sport competition effects. Real-world fan loyalty and sponsorship allocation may exhibit nonlinear dynamics not fully captured here.

Finally, long-term projections of player development and injury risk are probabilistic and subject to uncertainty. Although risk-adjusted mechanisms are embedded in the framework, forecasting error inevitably increases over extended horizons.

Overall, while the model provides a structured and quantitatively grounded decision-support system, practical implementation should combine analytical outputs with managerial judgment and updated real-time data.


\begin{letter}{Dear XXX:}

From simulation results of real-life situations, we find it takes a period of time for the inflow hot water to spread throughout the bathtub. During this process, the bath water continues transferring heat into air, bathtub and the person in bathtub. The difference between heat transfer capacity makes the temperature of various areas to be different. So that it is difficult to get an evenly maintained temperature throughout the bath water.

In order to enjoy a comfortable bath with even temperature of bath water and without wasting too much water, we propose the following suggestions.

\begin{itemize}
\item Adding hot water consistently
\item Using smaller bathtub if possible
\item Decreasing motions during bath
\item Using bubble bath additives
\item Arranging more faucets of inflow
\end{itemize}

\vspace{\parskip}

Sincerely yours,

Your friends

\end{letter}

\newpage

\begin{appendices}

\section{First appendix}

In addition, your report must include a letter to the Chief Financial Officer (CFO) of the Goodgrant Foundation, Mr. Alpha Chiang, that describes the optimal investment strategy, your modeling approach and major results, and a brief discussion of your proposed concept of a return-on-investment (ROI). This letter should be no more than two pages in length.

Here are simulation programmes we used in our model as follow.\\

\textbf{\textcolor[rgb]{0.98,0.00,0.00}{Input matlab source:}}
\lstinputlisting[language=Matlab]{./code/mcmthesis-matlab1.m}

\section{Second appendix}

some more text \textcolor[rgb]{0.98,0.00,0.00}{\textbf{Input C++ source:}}
\lstinputlisting[language=C++]{./code/mcmthesis-sudoku.cpp}

\end{appendices}

\newpage
\newcounter{lastpage}
\setcounter{lastpage}{\value{page}}
\thispagestyle{empty} 

\section*{Report on Use of AI}

\begin{enumerate}
\item OpenAI ChatGPT (Nov 5, 2023 version, ChatGPT-4,) 
\begin{description}
\item[Query1:] <insert the exact wording you input into the AI tool> 
\item[Output:] <insert the complete output from the AI tool>
\end{description}
\item OpenAI Ernie (Nov 5, 2023 version, Ernie 4.0) 
\begin{description}
\item[Query1:] <insert the exact wording of any subsequent input into the AI tool> 
\item[Output:] <insert the complete output from the second query>
\end{description}
\item Github CoPilot (Feb 3, 2024 version) 
\begin{description}
\item[Query1:] <insert the exact wording you input into the AI tool> 
\item[Output:] <insert the complete output from the AI tool>
\end{description}
\item Google Bard (Feb 2, 2024 version) 
\begin{description}
\item[Query1:] <insert the exact wording of your query> 
\item[Output:] <insert the complete output from the AI tool>
\end{description}
\end{enumerate}

% 重置页码
\clearpage
\setcounter{page}{\value{lastpage}}

\end{document}

%%
%% This work consists of these files mcmthesis.dtx,
%%                                   figures/ and
%%                                   code/,
%% and the derived files             mcmthesis.cls,
%%                                   mcmthesis-demo.tex,
%%                                   README,
%%                                   LICENSE,
%%                                   mcmthesis.pdf and
%%                                   mcmthesis-demo.pdf.
%%
%% End of file `mcmthesis-demo.tex'.
